\changecolours{ScoutNavy}
\section{O Modelo Chave-Valor e o Ecossistema Redis}

% ------------------------------------------------------------------------------
\twocolframe[title={O Paradigma Chave-Valor (Key-Value)},leftcol=7cm,rightcol=7cm]{%
  \alert{Simplicidade e Desempenho Extremo}\par
  \itemseps[topsepi=2mm,itemsepi=3mm,labelsep=2mm,leftmargini=4mm]
  \begin{itemize}
    \item Entre os 4 modelos NoSQL, o \textbf{Chave-Valor} é o mais primitivo e veloz.
    \item Cada dado é indexado por uma única chave opaca (\textit{Key}).
    \item O valor associado (\textit{Value}) é tratado como uma carga útil (\textit{payload}).
    \item \textbf{Complexidade de Busca:} Tempo constante $\mathcal{O}(1)$ via tabela hash interna.
  \end{itemize}
}{%
  \alert{Por que o Redis se Destaca no Mercado?}\par
  \itemseps[topsepi=2mm,itemsepi=3mm,labelsep=2mm,leftmargini=4mm]
  \begin{itemize}
    \item Criado por \textbf{Salvatore Sanfilippo (antirez)} em 2009.
    \item Sigla: \textbf{REmote DIctionary Server}.
    \item \textbf{In-Memory Engine:} Opera 100\% na memória RAM com latência sub-milissegundo ($< 1$ ms).
    \item Não armazena apenas strings brutas, mas \textbf{estruturas de dados ricas}.
  \end{itemize}
}

% ------------------------------------------------------------------------------
\twocolframe[title={Ambiente Windows: Por que Usar o WSL 2?},leftcol=7cm,rightcol=7cm]{%
  \alert{Natureza POSIX do Redis}\par
  \itemseps[topsepi=2mm,itemsepi=3mm,labelsep=2mm,leftmargini=4mm]
  \begin{itemize}
    \item O Redis depende de primitivas POSIX (Linux/Unix), como \texttt{fork()} para persistência em segundo plano.
    \item Não há binário oficial para Windows nativo desde a versão 3.2.
    \item Versões não-oficiais para Windows apresentam instabilidade e faltam recursos modernos.
  \end{itemize}
}{%
  \alert{A Solução Definitiva: WSL 2}\par
  \itemseps[topsepi=2mm,itemsepi=3mm,labelsep=2mm,leftmargini=4mm]
  \begin{itemize}
    \item O \textbf{Windows Subsystem for Linux 2} roda um kernel Linux real sobre hipervisor leve.
    \item \textbf{Desempenho Idêntico ao Linux:} Sem degradação em chamadas de sistema.
    \item \textbf{Rede Transparente:} O Windows mapeia \texttt{localhost:6379} automaticamente para IDEs e Python.
  \end{itemize}
}

% ------------------------------------------------------------------------------
\begin{frame}[fragile]{Instalação e Execução do Redis no WSL 2}
  \begin{columns}[T]
    \begin{column}{0.48\textwidth}
      \alert{1. Instalação do Ambiente Linux}\par
      \itemseps[topsepi=1mm,itemsepi=1.5mm,labelsep=2mm,leftmargini=4mm]
      \begin{itemize}
        \item No PowerShell (Administrador):
      \end{itemize}
      \begin{lstlisting}[language=bash]
# Instala o Ubuntu no WSL 2:
wsl --install -d Ubuntu
      \end{lstlisting}
      \vspace{0.1cm}
      \alert{2. Pacote Oficial no Ubuntu}\par
      \begin{lstlisting}[language=bash]
# Atualiza repositorios e instala:
sudo apt update
sudo apt install redis-server -y
      \end{lstlisting}
    \end{column}
    \begin{column}{0.50\textwidth}
      \alert{3. Gerenciamento do Serviço}\par
      \begin{lstlisting}[language=bash]
# Iniciar o servidor Redis:
sudo service redis-server start

# Verificar se esta rodando:
sudo service redis-server status
      \end{lstlisting}
      \vspace{0.1cm}
      \alert{4. Teste de Conexão}\par
      \begin{lstlisting}[language=bash]
# Abrir console interativo:
redis-cli
127.0.0.1:6379> PING
PONG
      \end{lstlisting}
    \end{column}
  \end{columns}
\end{frame}

% ------------------------------------------------------------------------------
\begin{frame}[fragile]{Convenção de Chaves e Diagnóstico no redis-cli}
  \begin{columns}[T]
    \begin{column}{0.48\textwidth}
      \alert{Convenção de Namespaces}\par
      \itemseps[topsepi=1mm,itemsepi=1.5mm,labelsep=2mm,leftmargini=4mm]
      \small
      \begin{itemize}
        \item O Redis não possui schemas relacionais.
        \item Chaves são organizadas com dois-pontos (\texttt{:}):
        \item \textbf{Padrão:} \texttt{\scriptsize modulo:entidade:id:campo}
        \item \textbf{Exemplo:} \texttt{\scriptsize loja:prod:42:preco}
        \item Garante legibilidade e isolamento.
      \end{itemize}
    \end{column}
    \begin{column}{0.50\textwidth}
      \alert{Diagnóstico no redis-cli}\par
      \begin{lstlisting}[language=bash]
# Testar eco e latencia:
127.0.0.1:6379> PING "IFCE Taua"
"IFCE Taua"

# Diagnostico do servidor:
127.0.0.1:6379> INFO server
# Retorna versao, PID e uptime

# Contagem total de chaves:
127.0.0.1:6379> DBSIZE
(integer) 42
      \end{lstlisting}
    \end{column}
  \end{columns}
\end{frame}
